% ====================================================================
% Section: The load-distribution attractor that explains the convergence.
% Faithful to data/attractor_full.json (12 seeds, cosine + L1, scipy tests).
% Robust claim: cores converge to one distribution, tighter than controls.
% Fragile claim (tube/sp link): reported with explicit caution.
% ====================================================================
\section{Why They Converge: A Load-Distribution Attractor}
\label{sec:attractor}

Section~\ref{sec:convergence} established \emph{that} three independent physical
cores reach the same drop-rate performance. This section asks \emph{why}, and
the answer turns the convergence from a coincidence into a mechanism: the three
cores do not merely score the same---they route traffic through the same edges
in the same proportions. The topology imposes a load-distribution attractor, and
the congestion-aware physics settles into it almost exactly.

\subsection{Measuring the distribution, not the score}

For each topology we run five routers on an identical graph and demand: the three
physical cores (Newton, Archimedes, Pascal) and two blind controls---shortest
path (ignores congestion) and ECMP (hash spread). For each router we record its
per-edge utilization vector, averaged over the simulation, and compare these
vectors between routers. Two routers with similar vectors are pushing load
through the same edges in the same proportions: they have found the same
distribution. We use two independent similarity measures so the result cannot
hinge on one definition: cosine similarity of the utilization vectors, and one
minus the normalized $L_1$ distance between them (how much normalized load mass
would have to be moved to turn one distribution into the other). Twelve seeds per
topology.

\subsection{The cores converge to one distribution}

Table~\ref{tab:attractor} reports the results. The three physical cores produce
near-identical load distributions on every topology: cosine similarity
$0.992 \pm 0.002$, and $0.956 \pm 0.007$ on the $L_1$ measure. The agreement is
both very high and very stable---the small standard deviations show it holds
across all fifteen topologies, not on average.

\begin{table}[t]
\centering
\caption{Similarity of per-edge load distributions between routers (12 seeds).
``phys--phys'' is the mean over the three physical-core pairs; ``phys--short''
and ``phys--ecmp'' compare the cores to the blind controls. Both similarity
measures agree: the cores converge to one distribution, tighter than either
control reaches.}
\label{tab:attractor}
\small
\begin{tabular}{lcc}
\toprule
 & cosine & $1-L_1$ \\
\midrule
physics--physics  & $0.992 \pm 0.002$ & $0.956 \pm 0.007$ \\
physics--shortest & $0.948 \pm 0.009$ & $0.846 \pm 0.021$ \\
physics--ECMP     & $0.866 \pm 0.054$ & $0.757 \pm 0.056$ \\
\bottomrule
\end{tabular}
\end{table}

Crucially, this convergence is tighter than what the blind controls reach. A
paired Wilcoxon test across topologies confirms the physical cores agree with
each other significantly more than they agree with shortest-path routing
($p = 0.0007$, on both similarity measures). The interpretation is twofold.
First, the topology does most of the work: even ECMP, which ignores congestion
entirely, lands at cosine $0.87$---the graph and the demand already constrain
where load can go. This is the structural attractor in its plain form, the
sense in which ``all roads lead to the same few places''. Second, the
congestion-aware physics does something the blind routers do not: it refines
that rough distribution into a near-exact consensus ($0.99$), and it is this
last refinement, shared by all three mechanisms, that produces the matched
drop rates of Section~\ref{sec:convergence}.

In other words, the three physics tie not by accident but because they descend
into the same basin: a load distribution fixed largely by the topology and
sharpened identically by any congestion-aware rule, regardless of which
classical mechanism implements it.

\subsection{A caveat we will not paper over}

One natural follow-up does \emph{not} survive scrutiny cleanly, and we report it
as such. We asked whether the attractor is stronger where the topology offers
less room to manoeuvre---whether divergence between cores grows with the tube/sp
ratio. Under the cosine measure the correlation is strong and significant
($r = +0.82$, $p = 0.0002$); under the $L_1$ measure it is weak and not
significant ($r = +0.37$, $p = 0.18$). Because our two measures disagree here,
we decline to claim a tube/sp--divergence law. We report the cosine correlation
as suggestive and leave the question open. The convergence result itself does
not depend on it: that the cores reach one distribution, more tightly than blind
controls, holds on both measures and at high significance.
